\section{Instalación y configuración de PostgreSQL}

\subsection{Descarga de la imagen}

PostgreSQL se desplegó dentro de un contenedor Docker utilizando la imagen
\texttt{postgres:18}. La etiqueta fija permite identificar la versión principal
empleada durante la práctica.

\begin{lstlisting}[style=terminal]
$ sudo docker pull postgres:18
\end{lstlisting}

\subsection{Credencial y creación del contenedor}

Se generó una contraseña aleatoria y se almacenó en un archivo local con
permisos restrictivos. El valor no se incluyó en la bitácora ni en las
capturas.

El primer intento de publicación utilizó el puerto local 5432, pero dicho
puerto ya estaba ocupado por otro contenedor llamado \texttt{postgres}. Para
evitar modificar o detener ese servicio, el contenedor de esta práctica se
publicó en el puerto local 5433, conservando el puerto interno 5432 de
PostgreSQL:

\begin{lstlisting}[style=terminal]
$ sudo docker run -d \\
    --name postgres-practica01 \\
    --env-file /home/maik/practica01-capturas/postgres.env \\
    -p 127.0.0.1:5433:5432 \\
    -v postgres_practica01_data:/var/lib/postgresql \\
    postgres:18
\end{lstlisting}

La opción de publicación se interpreta como
\texttt{localhost:5433 -> contenedor:5432}. El volumen permite conservar los
datos aunque el contenedor se detenga.

\evidencia{03_evidencia.png}{Contenedor \texttt{postgres-practica01} activo en el puerto local 5433 y PostgreSQL listo para aceptar conexiones.}

\subsection{Comprobación con psql}

La consulta se ejecutó dentro del contenedor para comprobar la base activa, el
usuario y la versión del servidor:

\begin{lstlisting}[style=terminal]
$ sudo docker exec postgres-practica01 \\
    psql -U postgres -d postgres \\
    -c 'SELECT current_database(), current_user, version();'
\end{lstlisting}

El resultado confirmó PostgreSQL 18.6 sobre Debian y la disponibilidad de la
base \texttt{postgres}.
